\documentclass[a4paper,11pt]{ltjsarticle}

% ---------- パッケージ ----------
\usepackage{luatexja}
\usepackage{luatexja-fontspec}
\usepackage{amsmath,amssymb,amsthm,mathtools}
\usepackage{bm}
\usepackage{geometry}
\usepackage{graphicx}
\usepackage{hyperref}
\usepackage{cleveref}
\usepackage{enumitem}

% ---------- 余白 ----------
\geometry{margin=25mm}

% ---------- フォント ----------
\setmainjfont{Noto Serif CJK JP} % 日本語
\setmainfont{Times New Roman}    % 欧文

% ---------- 行間 ----------
\linespread{1.2}

% ---------- 定理環境 ----------
\theoremstyle{plain}
\newtheorem{theorem}{定理}[section]
\newtheorem{lemma}[theorem]{補題}
\newtheorem{proposition}[theorem]{命題}

\theoremstyle{definition}
\newtheorem{definition}[theorem]{定義}
\newtheorem{example}[theorem]{例}

\theoremstyle{remark}
\newtheorem{remark}[theorem]{注意}

% ---------- 数学記号 ----------
\newcommand{\R}{\mathbb{R}}
\newcommand{\Z}{\mathbb{Z}}
\newcommand{\Q}{\mathbb{Q}}

% ---------- タイトル ----------
\title{\bfseries レポートタイトル}
\author{所属・氏名}
\date{\today}

% ---------- 本文 ----------
\begin{document}

\maketitle

\begin{abstract}
本レポートでは，○○について考察する．
特に，△△の性質および□□との関係を明らかにすることを目的とする．
\end{abstract}

\section{はじめに}

本レポートでは，○○について扱う．
この問題は，数学や情報科学において重要な役割を果たす．

\section{準備}

まず，基本的な定義を与える．

\begin{definition}
集合 $X$ 上の写像 $f : X \to X$ が性質 $(P)$ を満たすとは，
任意の $x \in X$ に対して
\[
  f(x) = x
\]
が成り立つことをいう．
\end{definition}

\begin{lemma}
$V$ を $\R$ 上のベクトル空間とする．
$T \in \mathrm{End}(V)$ が対角化可能であるとき，
\[
  V = \bigoplus_{\lambda} E_{\lambda}
\]
が成り立つ．
\end{lemma}

\begin{proof}
定義より従う．
\end{proof}

\section{主結果}

以下が本レポートの主結果である．

\begin{theorem}
仮定として○○が成り立つとする．
このとき，次が成立する：
\begin{enumerate}[label=(\arabic*)]
  \item ○○が成り立つ
  \item △△が成立する
\end{enumerate}
\end{theorem}

\begin{proof}
証明は以下の通りである．

まず，○○を示す．
次に，△△を示す．

したがって，定理が成立する．
\end{proof}

\section{例}

\begin{example}
$G = \mathrm{SL}_2(\Z)$ は，上半平面に自然に作用する：
\[
  z \mapsto \frac{az+b}{cz+d}
\]
\end{example}

\section{まとめ}

本レポートでは，○○について考察した．
今後の課題として，□□の拡張が考えられる．

\begin{thebibliography}{9}

\bibitem{sample}
著者名，
\textit{書名}，
出版社，年．

\end{thebibliography}

\end{document}